\cvsection{Expériences professionnelles}

\begin{cventries}
  \cventry
    {Stagiaire ingénieur de recherche en IA}
    {INRIA}
    {Bordeaux, France}
    {Mars 2026 -- Septembre 2026}
    {
      \begin{cvitems}
        \item {Analyse du besoin des chercheurs et conception en autonomie d’une application permettant d’identifier et d’exploiter des articles scientifiques pertinents.}
        \item {Automatisation de la collecte de métadonnées d’articles scientifiques disponibles en ligne avec tâches cron et transmission des résultats par webhooks.}
        \item {Mise en œuvre d’un pipeline RAG couvrant le chunking, la génération d’embeddings, l’indexation dans ChromaDB et la recherche sémantique.}
        \item {Construction sans framework d’orchestration d’un agent IA multi-étapes avec API OpenAI, prompts, outils, état, mémoire et gestion du contexte.}
        \item {Développement d’un backend Python/FastAPI avec API REST documentées avec OpenAPI et persistance des documents et métadonnées dans MongoDB.}
        \item {Développement d’une interface React/Next.js permettant aux chercheurs d’interagir avec l’agent, ses outils et les documents retrouvés.}
        \item {Évaluation de LLM, notamment Qwen, Gemma et GLM, et de plusieurs quantifications selon la qualité, la latence, le TTFT et la taille du contexte.}
        \item {Définition du protocole d’évaluation, analyse des résultats et rédaction en cours d’un article scientifique présentant la méthode et ses limites.}
      \end{cvitems}
    }

  \cventry
    {Stagiaire ingénieur de recherche en apprentissage par renforcement}
    {Université Karunya}
    {Coimbatore, Inde}
    {Juin 2025 -- Septembre 2025}
    {
      \begin{cvitems}
        \item {Développement d’un modèle d’apprentissage par renforcement pour la navigation autonome d’un robot hospitalier, puis déploiement et tests sur Raspberry Pi.}
      \end{cvitems}
    }
\end{cventries}
